% Shared canonical-ready fragment derived from two_cross_exceptional_exclusion.tex.
% Input immediately after q444:generic-projection-section and before the
% existing Sources, independent checks, and reproducibility subsection.
% No preamble, document environment, or global macro definitions.
% Stage the companion proof/checker/receipt at the relative evidence path below.
\subsection{An eight-dimensional exceptional spatial class with no full solution}
\label{q444:twocross:section}

\subsubsection{Literal system and the explicitly selected class}
\label{q444:twocross:literal-system}

All vector spaces in this note retain their real structures. In particular,
$\mathfrak{su}(3)$ is the eight-real-dimensional space of traceless anti-Hermitian
complex $3\times3$ matrices, with positive inner product
$\langle X,Y\rangle=-\operatorname{tr}(XY)$. Retain the original matrices
\[
H=\begin{pmatrix}0&0&0\\0&-1/2&1/(2\sqrt2)\\
0&1/(2\sqrt2)&1/2\end{pmatrix},\quad
q=\frac{e_1}{2\sqrt2},\quad
A=\sqrt2\begin{pmatrix}0&-1&0\\1&0&0\\0&0&0\end{pmatrix},
\]
and the original scalar constants $m^2=\Gamma=1$, $\kappa=6$.
For $Y_{\rm tf}=Y-\operatorname{tr}(Y)I_3/3$, define
\[
\mathcal IX=([H,[H,X]]+qq^*X+Xqq^*)_{\rm tf},\qquad
\mathcal KX=-[Z_1,[Z_1,X]]-[Z_2,[Z_2,X]],
\]
\[
F=[A,[A,H]]=-2H+\operatorname{diag}(-2,1,1).
\]
The last identity is verified by the two displayed matrix products; also
$H^2=(3/8)\operatorname{diag}(0,1,1)$ and $q^*q=1/8$.
The indeterminate/value $t=g^2$ is the original positive gauge coupling,
not physical time. The original frequency satisfies $\omega\in\mathbb R
\setminus\{0\}$. The five retained equations are
\begin{align}
[C,[C,H]]-[Z_1,[Z_1,H]]-[Z_2,[Z_2,H]]&=F,\label{q444:twocross:eq:H}\\
(C^2-Z_1^2-Z_2^2)q&=-2q,\label{q444:twocross:eq:q}\\
(t\mathcal I+2\mathcal K)C&=4\omega[Z_1,Z_2],\label{q444:twocross:eq:G}\\
-\omega^2Z_1-2\omega[C,Z_2]+[C,[C,Z_1]]
+\mathcal KZ_1+(t/2)\mathcal IZ_1&=0,\label{q444:twocross:eq:S1}\\
-\omega^2Z_2+2\omega[C,Z_1]+[C,[C,Z_2]]
+\mathcal KZ_2+(t/2)\mathcal IZ_2&=0.\label{q444:twocross:eq:S2}
\end{align}
Multiplying the original Gauss equation by its retained positive $t$
gives \eqref{q444:twocross:eq:G}; division by that same $t$ is its exact inverse.

Put $P=qq^*/(q^*q)=\operatorname{diag}(1,0,0)$ and define
\[
\mathfrak p=\{X\in\mathfrak{su}(3):PXP=(I_3-P)X(I_3-P)=0\},\qquad
\mathfrak k=\{X\in\mathfrak{su}(3):[X,P]=0\}.
\]
Thus $\mathfrak p$ exchanges the matter line and its orthogonal complement, while
$\mathfrak k$ preserves them. Every $X\in\mathfrak p$ has the unique form
\[
X(v)=\begin{pmatrix}0&-v^*\\v&0_{2\times2}\end{pmatrix},
\qquad v\in\mathbb C^2,
\]
whose inverse takes the lower first column. This is a real linear
isomorphism $\mathbb C^2\to\mathfrak p$, so $\mathfrak p\times\mathfrak p$ has real dimension eight.
Block decomposition gives the orthogonal direct sum $\mathfrak{su}(3)=\mathfrak k\oplus\mathfrak p$,
with both summands of real dimension four.

\begin{theorem}[Two-cross spatial exclusion]\label{q444:twocross:thm:exclusion}
There is no solution of \eqref{q444:twocross:eq:H}--\eqref{q444:twocross:eq:S2} with both
$Z_1,Z_2\in\mathfrak p$. No hypothesis on $C$ is imposed in this statement.
Every complex coordinate of each of the two cross vectors is allowed,
and both signs of $\omega\ne0$ and every original $t>0$ are retained.
\end{theorem}

This is an explicitly restricted spatial class, not a claim that arbitrary
spatial pairs can be put into this class. The proof below also shows that
the entire class lies in the previously defined real algebraic exceptional
zero set. Consequently the result excludes an eight-dimensional linear
subspace inside that zero set, but does not decide existence in its
remaining part.

\subsubsection{Gauss and the complete scalar equations force the temporal form}
\label{q444:twocross:temporal-form}

Block multiplication gives
$[\mathfrak p,\mathfrak p]\subset\mathfrak k$, $[\mathfrak k,\mathfrak p]\subset\mathfrak p$.
Because $H$ and $qq^*$ are block diagonal, $\mathcal I$ preserves
both summands. Direct multiplication using $H^2$ and $q^*q$ gives
\begin{equation}
\mathcal IX(v)=\left(\frac38+\frac18\right)X(v)=\frac12X(v).
\label{q444:twocross:eq:Ip}
\end{equation}
In this computation $HX(v)H=0$, the two $H^2$ products sum to
$(3/8)X(v)$, the two $qq^*$ products sum to $X(v)/8$, and every
term is traceless. If $Z_1,Z_2\in\mathfrak p$, their double commutators
therefore make $\mathcal K$ preserve both summands as well.
Cyclicity gives, for every $Y\in\mathfrak{su}(3)$,
\[
\langle Y,\mathcal KY\rangle=\sum_{j=1}^2\|[Z_j,Y]\|^2\ge0.
\]
Project \eqref{q444:twocross:eq:G} onto $\mathfrak p$. Its right side lies in $\mathfrak k$, hence
\[
\left(\frac t2\operatorname{id}_{\mathfrak p}+2\mathcal K|_{\mathfrak p}\right)C_{\mathfrak p}=0.
\]
Pairing with $C_{\mathfrak p}$ gives a sum whose first term is
$(t/2)\|C_{\mathfrak p}\|^2$ and whose other terms are nonnegative.
Since $t>0$, $C_{\mathfrak p}=0$. Thus Gauss forces $C\in\mathfrak k$; this was
not an additional restriction on the temporal unknown.

Here are all frame changes, together with their inverses. Put
\[
x=\frac12,\quad d_0=\frac1{2\sqrt2},\quad r=\frac{\sqrt6}{4},
\quad \delta=\frac1{2r}=\sqrt{\frac23},
\]
\[
S=\frac1{\sqrt{2r(r+x)}}
\begin{pmatrix}d_0&-(r+x)\\r+x&d_0\end{pmatrix},\qquad
U=\operatorname{diag}(1,S).
\]
The relation $r^2=x^2+d_0^2$ shows that the numerator columns have
squared length $2r(r+x)$, are orthogonal, and have determinant
$2r(r+x)>0$. Multiplication by the lower block of $H$ gives their
eigenvalues $r,-r$ in that order. Hence
\[
U\in SO(3),\qquad U^*HU=\operatorname{diag}(0,r,-r),\qquad U^*q=q.
\]
For every matrix $Y$ the change is $Y\mapsto U^*YU$, with inverse
$Y'\mapsto UY'U^*$. It is applied to $H,F,C,Z_1,Z_2$ simultaneously.
It commutes with products, brackets, adjoints and trace projection,
and fixes $q$. Thus it intertwines $\mathcal I,\mathcal K$ and each
of the five full residuals exactly. It preserves $\mathfrak p,\mathfrak k$ because
it commutes with $P$. No physical constant is rescaled.

In this first frame write the full temporal block matrix as
\[
C_d=\begin{pmatrix}
ia&0&0\\0&i(h-a)/2&z\\0&-\bar z&-i(a+h)/2
\end{pmatrix},\qquad a,h\in\mathbb R,\quad z\in\mathbb C.
\]
Write the spatial matrices as $X(v),X(w)$ and put
$M=vv^*+ww^*$ in this frame only for the following calculation.
The vector scalar residual is exactly
\[
(C_d^2-X(v)^2-X(w)^2)q+2q
=(-a^2+\operatorname{tr}M+2)q.
\]
The complete nonzero block of the adjoint scalar residual, including
its upper entry, is
\[
\begin{pmatrix}
2-2r(M_{11}-M_{22})&0&0\\
0&-4r|z|^2+2rM_{11}-1+2r&-2irhz\\
0&2irh\bar z&4r|z|^2-2rM_{22}-1-2r
\end{pmatrix}.
\]
For example, if the lower Hermitian block $B=-i(C_d)_{\rm lower}$
and $D=\operatorname{diag}(r,-r)$ are used, its lower block is
$-[B,[B,D]]+\{M,D\}-(I_2-2D)$; multiplying these two-by-two
matrices gives every entry just displayed.
The two scalar equations are therefore equivalent to
\begin{equation}
a^2=2+\operatorname{tr}M,\qquad
2|z|^2-M_{11}=1-\delta,\qquad
2|z|^2-M_{22}=1+\delta,\qquad hz=0.
\label{q444:twocross:eq:first-scalar}
\end{equation}
The upper entry is minus the sum of the two lower diagonal entries,
so no scalar entry is left unchecked. Adding the middle two equations
and comparing with the first gives $a^2=4|z|^2$. Since
$M_{22}\ge0$ and $1+\delta>0$, $z\ne0$. Consequently
\begin{equation}
h=0,\qquad a\ne0,\qquad |z|=|a|/2.
\label{q444:twocross:eq:temporal-forced}
\end{equation}

Set $e=a$, retaining either sign, and $\gamma=-2iz/e$, so
$|\gamma|=1$. Choose either square root $\zeta$ of $-\bar\gamma$ and put
\[
W=\frac{\zeta}{\sqrt2}
\begin{pmatrix}\gamma&\gamma\\1&-1\end{pmatrix},\qquad
V=\operatorname{diag}(1,W).
\]
Then $|\zeta|=1$, $W^*W=I_2$ and
$\det W=-\zeta^2\gamma=\bar\gamma\gamma=1$.
The lower Hermitian block of $C_d$ is
\[
\frac e2\begin{pmatrix}-1&\gamma\\\bar\gamma&-1\end{pmatrix}.
\]
Its products with $(\gamma,1)^{\mathsf T}$ and
$(\gamma,-1)^{\mathsf T}$ have eigenvalues $0,-e$, respectively.
Those same columns give
$W^*\operatorname{diag}(r,-r)W=r\left(\begin{smallmatrix}0&1\\1&0\end{smallmatrix}\right)$.
Thus the second exact change $Y_d\mapsto V^*Y_dV$ gives
\begin{equation}
C'=ie\operatorname{diag}(1,0,-1),\quad
H'=\begin{pmatrix}0&0&0\\0&0&r\\0&r&0\end{pmatrix},\quad q'=q,
\quad F'=-2H'+\operatorname{diag}(-2,1,1).
\label{q444:twocross:eq:second-frame}
\end{equation}
The inverse for every original matrix is, explicitly,
\begin{equation}
Y=UVY'V^*U^*.
\label{q444:twocross:eq:original-inverse}
\end{equation}
For spatial vectors the forward change is the corresponding lower
unitary action $v\mapsto W^*v$ and its inverse is $v'\mapsto Wv'$.
The unitary parameters are retained; the assertion is not that $H$ can
be changed while the other data are left fixed. As with $U$, this
simultaneous conjugation intertwines all five equations. In the remainder
of the proof primes are omitted only to ease typography; all matrices
are the exact objects in \eqref{q444:twocross:eq:second-frame} and have inverse
\eqref{q444:twocross:eq:original-inverse}.

\subsubsection{Every gauge entry in the classified frame}
\label{q444:twocross:full-gauge-entries}

Write $Z_1=X(v)$, $Z_2=X(w)$ with arbitrary $v,w\in\mathbb C^2$,
and define, afresh in the second frame,
\[
V_0=v^*v,\quad W_0=w^*w,\quad p_0=v^*w=\alpha_0+ib,
\quad M=vv^*+ww^*,\quad N=i(vw^*-wv^*).
\]
Both $M$ and $N$ are Hermitian; $M\succeq0$. Put
\[
u=M_{11},\quad v_0=M_{22},\quad m=M_{12},\quad
d_M=\det M=|v_1w_2-v_2w_1|^2\ge0.
\]
The determinant identity follows by expanding the four products in
$M_{11}M_{22}-M_{12}M_{21}$ and collecting the squared modulus.
The symbol $v_0$ here is a real diagonal entry, not the vector $v$.
All eight real components of the two vectors remain independent.

Multiplication gives the full scalar residuals in this frame:
\[
Q_H=\begin{pmatrix}
2-4r\operatorname{Re}m&0&0\\
0&2r\operatorname{Re}m-1&r(\operatorname{tr}M-e^2+2)\\
0&r(\operatorname{tr}M-e^2+2)&2r\operatorname{Re}m-1
\end{pmatrix},
\quad Q_q=(\operatorname{tr}M-e^2+2)q.
\]
Consequently scalar compatibility is exactly
\begin{equation}
e^2=2+\operatorname{tr}M,\qquad \operatorname{Re}m=\delta>0.
\label{q444:twocross:eq:second-scalar}
\end{equation}
In particular $m\ne0$. No phase of $m,v,w$ has been set to a chosen value.

Set
\[
Q=\operatorname{diag}(1,2),\qquad
E=\operatorname{diag}(2/3,-5/6).
\]
Direct multiplication, keeping $r^2=3/8$ and $q^*q=1/8$, gives
\begin{equation}
\mathcal IC=ie\operatorname{diag}(1/6,2/3,-5/6),\qquad
\mathcal IX(v)=\tfrac12X(v).
\label{q444:twocross:eq:inertia-classified}
\end{equation}
For the first formula the two contributions, respectively from $H$
and from the trace-free $q$ term, are
\[
 ie\operatorname{diag}(0,3/4,-3/4),\qquad
 ie\operatorname{diag}(1/6,-1/12,-1/12).
\]
The second formula is \eqref{q444:twocross:eq:Ip}, or follows from
$(H')^2=(3/8)\operatorname{diag}(0,1,1)$ directly.
The entire spatial bracket and the entire $\mathcal KC$ are
\[
[Z_1,Z_2]=\begin{pmatrix}-2ib&0\\0&iN\end{pmatrix},\qquad
\mathcal KC=\begin{pmatrix}2ie\operatorname{tr}(QM)&0\\
0&-ie(QM+MQ)\end{pmatrix}.
\]
For example $[C,X(v)]=X(-ieQv)$, and the two double cross products
give the displayed anticommutator. Thus the full Gauss equation is
equivalent to the upper entry and lower Hermitian block
\begin{equation}
\frac{et}{6}+4e\operatorname{tr}(QM)+8\omega b=0,\qquad
etE-2e(QM+MQ)-4\omega N=0.
\label{q444:twocross:eq:gauss-block}
\end{equation}
The upper entry is the negative lower trace, since
$\operatorname{tr}N=2b$. No off-block entries occur. The entries needed
below are therefore exact consequences of the full Gauss equation:
\begin{equation}
b=-\frac e{2\omega}\left(u+2v_0+\frac t{24}\right),\qquad
N_{12}=-\frac{3e}{2\omega}m,\qquad
N_{22}=-\frac e\omega\left(2v_0+\frac{5t}{24}\right).
\label{q444:twocross:eq:gauss-needed}
\end{equation}
All divisions use the original $\omega\ne0$.

Put
\[
D=\left(\frac t4-\omega^2\right)I_2-e^2Q^2.
\]
The complete spatial residuals are $X(s_1),X(s_2)$, respectively, where
\begin{align}
s_1&=(D+W_0I_2)v+
 \left(-\alpha_0I_2+i(2\omega eQ-3bI_2)\right)w,\label{q444:twocross:eq:s1}\\
s_2&=(D+V_0I_2)w+
 \left(-\alpha_0I_2-i(2\omega eQ-3bI_2)\right)v.\label{q444:twocross:eq:s2}
\end{align}
To verify every term, $[C,X(v)]=X(-ieQv)$ gives the two gyroscopic
terms and the term $-e^2Q^2$. The inertia gives $tI_2/4$.
For the remaining double commutators, block multiplication gives
\[
\mathcal KX(v)=X\bigl(W_0v+(\bar p_0-2p_0)w\bigr),\qquad
\mathcal KX(w)=X\bigl(V_0w+(p_0-2\bar p_0)v\bigr).
\]
Since $\bar p_0-2p_0=-\alpha_0-3ib$ and
$p_0-2\bar p_0=-\alpha_0+3ib$, these are exactly
\eqref{q444:twocross:eq:s1}--\eqref{q444:twocross:eq:s2}. Their upper rows are the negative
adjoints of their lower first columns, and all other entries are zero.
Thus $s_1=s_2=0$ is equivalent to the two full spatial equations.

We now use a consequence of those equations; no converse is needed.
Multiplying \eqref{q444:twocross:eq:s1} by $v^*$, \eqref{q444:twocross:eq:s2} by $w^*$ and adding
gives the exact matrix identity
\begin{equation}
s_1v^*+s_2w^*
=DM+d_MI_2+(4bI_2-2\omega eQ)N.
\label{q444:twocross:eq:outer-product}
\end{equation}
Indeed $\operatorname{tr}(M)M-M^2=d_MI_2$ follows by multiplying
the two-by-two matrix $M$. Its off-diagonal entries are zero, and its
diagonal entries are both $M_{11}M_{22}-M_{12}M_{21}$.
The cubic terms on the left of \eqref{q444:twocross:eq:outer-product} are
\[
W_0vv^*+V_0ww^*-\alpha_0(vw^*+wv^*)+3ib(vw^*-wv^*).
\]
Those in $\operatorname{tr}(M)M-M^2$ have $-ib$ instead of $3ib$;
their difference is $4bN$. The gyroscopic terms sum to
$-2\omega eQN$. These expansions prove \eqref{q444:twocross:eq:outer-product}
entry by entry, without a reality or orthogonality restriction on $v,w$.

\begin{proof}[Proof of Theorem~\ref{q444:twocross:thm:exclusion}]
Gauss and the scalar equations have already forced the exact temporal
form and \eqref{q444:twocross:eq:second-scalar}. In particular $e\ne0$, $m\ne0$.
For a full solution, \eqref{q444:twocross:eq:outer-product} is zero. Set
\[
\chi=\frac{3e}{2\omega},\qquad B_2=4b-4\omega e,
\quad D_2=\frac t4-\omega^2-4e^2.
\]
Its $(2,1)$ entry is
$D_2\bar m+B_2N_{21}=\bar m(D_2-\chi B_2)=0$ by Gauss.
Since $m\ne0$, it follows that $D_2=\chi B_2$.
Its $(2,2)$ entry is therefore
\[
0=d_M+B_2(\chi v_0+N_{22}).
\]
Using exactly \eqref{q444:twocross:eq:gauss-needed} gives
\[
\chi v_0+N_{22}=-\frac e{2\omega}\left(v_0+\frac{5t}{12}\right),
\qquad
B_2=-\frac{2e}{\omega}
 \left(u+2v_0+\frac t{24}+2\omega^2\right).
\]
Consequently every proposed full solution in this class would satisfy
\begin{equation}
0=d_M+\frac{e^2}{\omega^2}
\left(u+2v_0+\frac t{24}+2\omega^2\right)
\left(v_0+\frac{5t}{12}\right).
\label{q444:twocross:eq:positive-contradiction}
\end{equation}
But $u,v_0\ge0$, $d_M\ge0$, $t>0$, $\omega^2>0$ and $e^2>0$.
Both parentheses are strictly positive. The second term on the right
is strictly positive for either sign of $e$ and either sign of $\omega$,
contradicting \eqref{q444:twocross:eq:positive-contradiction}. This proves the theorem
without selecting a frequency or deleting any cross-vector phase.
\end{proof}

\begin{corollary}[A further necessary condition on every full solution]
Every solution of the unrestricted five equations satisfies
\[
\sum_{j=1}^2\left(
\|PZ_jP\|_{\rm HS}^2+
\|(I_3-P)Z_j(I_3-P)\|_{\rm HS}^2\right)>0.
\]
\end{corollary}
\begin{proof}
Each summand is a squared Hilbert--Schmidt norm and is nonnegative.
The sum vanishes if and only if both spatial matrices lie in $\mathfrak p$,
which is excluded by Theorem~\ref{q444:twocross:thm:exclusion}.
\end{proof}

The scalar-compatible part of the class is nonempty even with a nonzero
spatial bracket. This is not a full solution: in the frame
\eqref{q444:twocross:eq:second-frame}, choose the exact values
\[
v=\sqrt\delta(1,1)^{\mathsf T},\qquad w=(i,0)^{\mathsf T},
\qquad e=\sqrt{3+2\delta}.
\]
Then $\operatorname{tr}M=1+2\delta$, $M_{12}=\delta$, and
$[X(v),X(w)]_{11}=-2i\sqrt\delta\ne0$.
Thus both full scalar equations hold by their displayed residuals.
To embed this example into the literal frame, use the fixed
$W=i\left(\begin{smallmatrix}1&1\\1&-1\end{smallmatrix}\right)/\sqrt2$
and \eqref{q444:twocross:eq:original-inverse}. The choice of the component $i$ is a
specified example, not a normalization of the general class. The theorem
proves that no retained $t>0,\omega\ne0$ can complete this example to all
five equations.

\subsubsection{Exact inclusion of the selected class in the exceptional locus}
\label{q444:twocross:exceptional-inclusion-section}

We give the determinant map explicitly to specify the comparison with the
full sixteen-coordinate projection, including its degree conventions.
In this section the symbols $H,q,F$ again denote their literal original
matrices in \eqref{q444:twocross:eq:H}--\eqref{q444:twocross:eq:S2},
identical to the earlier retained data of
\eqref{q444:cx:eq:phi}--\eqref{q444:cx:eq:sp2}. For any such matrix write
$\widetilde Y=U^*YU$ when using its first, $H$-diagonal-frame value.
The inverse remains $Y=U\widetilde YU^*$. Put
\[
J=iH,\quad j_0=\langle J,J\rangle=\frac34,\quad
\mathscr W=J^\perp,\quad
\Pi X=X-\frac{\langle J,X\rangle}{j_0}J,\quad \mathscr B=[Z_1,Z_2],
\]
\[
\eta=\langle J,\mathcal KJ\rangle,\quad w_J=\Pi\mathcal KJ,
\quad\beta=\langle J,\mathscr B\rangle,\quad
\mathcal I_7=\mathcal I|_{\mathscr W},
\]
\[
\mathcal S_7=\Pi\mathcal K|_{\mathscr W}
-\frac{w_J\otimes w_J}{\eta},\qquad
f_7=\Pi\mathscr B-\frac\beta\eta w_J.
\]
Here $(w_J\otimes w_J)X=w_J\langle w_J,X\rangle$ and the original length of
$J$ is retained. On $\eta>0$ define
\[
\mathscr M_7(t)=t\mathcal I_7+2\mathcal S_7,\qquad
D_7=\det_{\mathscr W}\mathscr M_7,\qquad
V_7=4\operatorname{adj}_{\mathscr W}(\mathscr M_7)f_7,
\]
\[
N_7=V_7+\frac{2\beta D_7-\langle w_J,V_7\rangle}{\eta}J,
\qquad Q_7=\|V_7q\|^2,
\quad\alpha=\frac14+\|Z_1q\|^2+\|Z_2q\|^2,
\]
\[
P_H(t)=\alpha[N_7,[N_7,H]]
-Q_7\left(\sum_{j=1}^2[Z_j,[Z_j,H]]+F\right).
\]
Let $F_\xi=\eta^{16}(\widetilde P_H)_{11}$ and
$G_\xi=\eta^{16}(\widetilde P_H)_{22}$ denote the exact polynomials used for
the original sixteen-coordinate exceptional set, with inverse division
by $\eta^{16}$ on $\eta>0$. The matrix $P_H$ in its definition is
computed with the literal original $H,q,F$; $\widetilde P_H=U^*P_HU$
is then its exact conjugate, not a second change of already-conjugated data.
They extend polynomially in the sixteen real spatial coordinates $\xi$.
For clarity, their polynomial extension follows from the definitions:
$\eta^7D_7$, $\eta^7V_7$, $\eta^8N_7$ and $\eta^{14}Q_7$ are polynomial,
by determinant and minor expansions. Hence $\eta^{16}P_H$ is polynomial.
All other denominators are fixed nonzero constants. Their two selected
diagonal entries are real because $P_H$ is Hermitian for real $t$.
The universal degree bounds are $\deg V_7\le6$ and
$N_7=V_7+s_7J$ with $\deg s_7\le7$. Since
$[\widetilde J,\widetilde H]=0$ and every diagonal entry of
$[\widetilde J,Y]$ is zero, the diagonal entries of
$[\widetilde N_7,[\widetilde N_7,\widetilde H]]$ have degree at most
twelve; also $\deg Q_7\le12$. Thus $F_\xi,G_\xi$ have degree at most
twelve in $t$, including every specialization where their degrees drop.

The exceptional polynomial is the determinant $\mathscr R(\xi)$ of the
fixed map
\begin{equation}
\mathbb R[t]_{<12}\oplus\mathbb R[t]_{<12}\longrightarrow
\mathbb R[t]_{<24},\qquad (A_0,B_0)\longmapsto A_0F_\xi+B_0G_\xi,
\label{q444:twocross:eq:fixed-map}
\end{equation}
in increasing-power bases, with the first summand before the second.
Theorem~\ref{q444:generic-projection-obstruction} proves this polynomial
is nonzero on the ambient sixteen-dimensional carrier and that all full
solutions project to its zero set. The following inclusion is proved directly, independently
of whether any point of that zero set admits a solution.

\begin{proposition}\label{q444:twocross:prop:exceptional-inclusion}
The exact linear embedding $\mathfrak p\times\mathfrak p\hookrightarrow\mathfrak{su}(3)^2$ has image
inside $\{\mathscr R=0\}$. On every nonzero pair in this class the two
diagonal Gauss-eliminated polynomials have the common factor $D_4(t)^2$,
where
\[
D_4(t)=\det_{\mathfrak p}\left(\frac t2\operatorname{id}_{\mathfrak p}+2\mathcal K|_{\mathfrak p}\right)
\]
has degree exactly four and leading coefficient $1/16$.
\end{proposition}
\begin{proof}
The first-frame $\widetilde J=i\operatorname{diag}(0,r,-r)$ satisfies
$[\widetilde J,X(v)]=X(i\operatorname{diag}(r,-r)v)$.
The change by $U$ preserves $\mathfrak p$ and the inner product.
Thus, for the selected class,
\[
\eta=\sum_j\|[Z_j,J]\|^2
=r^2(\|Z_1\|^2+\|Z_2\|^2).
\]
It is positive except at the zero pair. This also shows why this class is
compatible with the stabilizer decomposition: $-\operatorname{ad}_J^2$
is exactly $r^2\operatorname{id}$ on $\mathfrak p$.

Put $\mathfrak k_0=\mathfrak k\cap J^\perp$. The orthogonal decomposition is
$\mathscr W=\mathfrak k_0\oplus\mathfrak p$, with dimensions three and four.
The already-proved block preservation of $\mathcal I,\mathcal K$,
and $J\in\mathfrak k$, give $w_J,f_7\in\mathfrak k_0$.
Consequently $\mathcal S_7$ preserves both summands and equals
$\mathcal K$ on $\mathfrak p$. Thus
\[
\mathscr M_7=\begin{pmatrix}\mathscr M_3&0\\0&\mathscr M_4\end{pmatrix},
\qquad \mathscr M_4=\frac t2\operatorname{id}_{\mathfrak p}+2\mathcal K|_{\mathfrak p},
\quad D_7=D_3D_4.
\]
Its four-dimensional determinant has degree four and leading coefficient
$(1/2)^4=1/16$. Moreover $\mathscr M_4$ is positive for $t>0$.
For completeness, $\mathcal I$ has kernel $\mathbb RJ$: its quadratic
form is $\operatorname{tr}[X,H]^2+2\|Xq\|^2$; a zero commutes with the
three-distinct-eigenvalue $H$ and annihilates $q$, so has diagonal
entries $0,is,-is$. The same formula proves positivity on $\mathscr W$.
Completing the square in $\langle X+sJ,\mathcal K(X+sJ)\rangle$ gives
\[
\langle X,\mathcal S_7X\rangle
+\eta\left(s+\frac{\langle w_J,X\rangle}{\eta}\right)^2\ge0,
\]
so choosing the displayed square to be zero proves
$\mathcal S_7\succeq0$. Thus $\mathscr M_7$ and its three-dimensional
block are positive for every $t>0$.

Define $V_3=4\operatorname{adj}(\mathscr M_3)f_7\in\mathfrak k_0$.
The block inverse for $t>0$ gives
\[
V_7=D_4V_3,\qquad
N_7=D_4N_3,\qquad
N_3=V_3+\frac{2\beta D_3-\langle w_J,V_3\rangle}\eta J.
\]
These are polynomial identities in $t$: the two sides agree for every
positive $t$, while a nonzero polynomial has only finitely many roots,
as follows by repeated division by its linear root factors.
Therefore $Q_7=D_4^2Q_3$, where $Q_3=\|V_3q\|^2$, and
\begin{equation}
P_H=D_4^2P_{H,3},\qquad
P_{H,3}=\alpha[N_3,[N_3,H]]
-Q_3\left(\sum_j[Z_j,[Z_j,H]]+F\right).
\label{q444:twocross:eq:factor}
\end{equation}
The three-dimensional minors give $\deg V_3\le2$, $\deg D_3\le3$.
Writing $N_3=V_3+s_3J$, one has $[J,H]=0$. In the diagonal frame,
$[J,Y]_{jj}=0$ for every matrix $Y$, by subtracting the identical
diagonal products. Hence the selected diagonal entries of
$[N_3,[N_3,H]]=[V_3,[V_3,H]]+s_3[J,[V_3,H]]$ have degree at most four.
Also $\deg Q_3\le4$. Thus at each fixed pair with $\eta>0$,
\[
F_\xi=D_4^2F_3,\qquad G_\xi=D_4^2G_3,\qquad
\deg F_3,\deg G_3\le4,
\]
where $F_3,G_3$ retain the scalar multiplier $\eta^{16}$.
If they are not both zero, $(A_0,B_0)=(G_3,-F_3)$ is a nonzero
kernel element of \eqref{q444:twocross:eq:fixed-map}, with both degrees less than
twelve. If both are zero, $(1,0)$ is such a kernel element instead.
The determinant therefore vanishes for every nonzero pair in the class,
including every degree-drop case. At the zero pair it vanishes by
continuity of the original polynomial $\mathscr R$ along this linear
subspace. This proves the entire asserted inclusion without using a
converse resultant test or interpreting $\eta=0$ as an inverse solution.
\end{proof}

\subsubsection{An exact certificate retaining nonzero residuals}
\label{q444:twocross:quantitative-section}

The exclusion has a quantitative extension. It also gives a polynomial
identity before the Gauss or spatial residuals are set to zero. Use the
entire second-frame variables above, with every entry of $v,w$ retained.
More precisely, return here to the matrices and inverse in
\eqref{q444:twocross:eq:second-frame}--\eqref{q444:twocross:eq:original-inverse},
not the literal-frame $H$ used in the preceding determinant comparison.
All of $M,N,m,u,v_0,b,s_1,s_2$ in this subsection are computed in
that same second frame as in \eqref{q444:twocross:eq:s1}--\eqref{q444:twocross:eq:s2}.
Define the full spatial residual matrix and the Hermitian Gauss components
\[
 \mathcal A=s_1v^*+s_2w^*,\qquad
 \gamma_0=et/6+4e(u+2v_0)+8\omega b,
 \qquad \mathcal G=etE-2e(QM+MQ)-4\omega N.
\]
Here $\mathcal A$ is not the original generator $A$: its calligraphic
symbol denotes exactly this two-by-two residual matrix. The full Gauss
residual is $i\operatorname{diag}(\gamma_0,\mathcal G)$, with the
second block of size two. Set $B_2=4b-4\omega e$ and
\[
 \mathcal P=d_M+\frac{e^2}{\omega^2}
       (u+2v_0+t/24+2\omega^2)(v_0+5t/12).
\]
This expression is defined for the original nonzero $\omega$; the
identity below has no such denominator after multiplication.

\begin{proposition}[Polynomial residual certificate and strict error bound]
\label{q444:twocross:quantitative}
For every $v,w\in\mathbb C^2$ and real $e,t,\omega$ in the displayed
second-frame matrices, the following polynomial identity holds:
\begin{align}
 &4\omega^2(\bar m\mathcal A_{22}-v_0\mathcal A_{21})
   +\omega B_2(\bar m\mathcal G_{22}-v_0\mathcal G_{21})\notag\\
 &\hspace{8mm}+e(v_0+5t/12)\bar m\gamma_0\notag\\
 &=4\bar m\bigl[\omega^2d_M+
      e^2(u+2v_0+t/24+2\omega^2)(v_0+5t/12)\bigr].
 \label{q444:twocross:polynomial-certificate}
\end{align}
No scalar, Gauss or spatial equation is assumed in this identity.

Now restrict to the exactly specified set of original tuples with
$Z_1,Z_2\in\mathfrak p$, $t>0$, $\omega\ne0$, for which both
scalar equations and Gauss hold. Let $R_1,R_2\in\mathfrak{su}(3)$
be their two full spatial residuals. Then
\begin{equation}
 \bigl(\|R_1\|_{\mathrm{HS}}^2+
             \|R_2\|_{\mathrm{HS}}^2\bigr)^{1/2}
 \geq
 \frac{\sqrt2\,\mathcal P}
 {\sqrt{\operatorname{tr}M\,(1+v_0^2/|m|^2)}}
 >\frac{5\sqrt2\,e^2t}
 {6\sqrt{\operatorname{tr}M\,(1+v_0^2/|m|^2)}}.
 \label{q444:twocross:residual-lower-bound}
\end{equation}
The full matrix norms are preserved by the two exact unitary frame maps.
In particular, on any such set with $\operatorname{tr}M\leq L$ and
$t\geq t_0>0$, where $L>0$, the strictly positive uniform bound is
\begin{equation}
 \bigl(\|R_1\|_{\mathrm{HS}}^2+
             \|R_2\|_{\mathrm{HS}}^2\bigr)^{1/2}
 >\frac{5\sqrt2\,t_0}
 {3\sqrt{L(1+L^2/\delta^2)}},\qquad \delta=1/(2r).
 \label{q444:twocross:uniform-bound}
\end{equation}
The complete scalar--Gauss solution set is parameterized in
Theorems~\ref{sgd:parametrization} and~\ref{sgd:literal-addon:quotient},
and Corollary~\ref{sgd:every-coupling} proves it is nonempty at every
original $t=g^2>0$. No sequence in its displayed bounded portion can
make both spatial residuals tend to zero.
\end{proposition}

\begin{proof}
The complete residual formulas above give
\[
 \mathcal A_{21}=D_2\bar m+B_2N_{21},\qquad
 \mathcal A_{22}=D_2v_0+d_M+B_2N_{22},\quad
 D_2=t/4-\omega^2-4e^2,
\]
and, without any equation imposed,
\[
 \mathcal G_{21}=-6e\bar m-4\omega N_{21},\qquad
 \mathcal G_{22}=-5et/6-8ev_0-4\omega N_{22}.
\]
Subtracting the two indicated spatial entries cancels $D_2$ exactly:
\[
 \bar m\mathcal A_{22}-v_0\mathcal A_{21}
 =\bar m d_M+B_2(\bar mN_{22}-v_0N_{21}).
\]
Multiplication of the indicated Gauss entries gives
\[
 \bar m\mathcal G_{22}-v_0\mathcal G_{21}
 =-e(2v_0+5t/6)\bar m
      -4\omega(\bar mN_{22}-v_0N_{21}).
\]
These cancel the two $N$ entries on the left of
\eqref{q444:twocross:polynomial-certificate}. The remaining factor is
\[
 \gamma_0-2\omega B_2
 =4e(u+2v_0+t/24+2\omega^2).
\]
It yields exactly its stated right side. Each displayed step is a
polynomial identity, including at $m=0$ or $\omega=0$; no inverse or
solution is assigned at those values.

On the stated scalar--Gauss zero set the preceding proof forces the
second-frame form, $e^2=2+\operatorname{tr}M$ and
$\operatorname{Re}m=\delta>0$. Hence $m\ne0$ and
$\operatorname{tr}M>0$. Setting $\gamma_0=\mathcal G=0$ in the
polynomial identity and dividing by its retained nonzero factors gives
\[
 \mathcal A_{22}-\frac{v_0}{\bar m}\mathcal A_{21}
 =\mathcal P>0.
\]
Cauchy--Schwarz on these two entries implies
\[
 \mathcal P\leq
  \sqrt{1+v_0^2/|m|^2}\,\|\mathcal A\|_{\mathrm{HS}}.
\]
Let $S=(s_1\ s_2)$ and $V=(v\ w)$ be the two column matrices.
Then $\mathcal A=SV^*$. Applying Cauchy--Schwarz to each entry and
summing proves $\|SV^*\|_{\mathrm{HS}}^2
\leq\|S\|_{\mathrm{HS}}^2\|V\|_{\mathrm{HS}}^2$.
The right factors are $\|s_1\|^2+\|s_2\|^2$ and
$\operatorname{tr}M$, respectively. Directly from the two off-diagonal
blocks of $X(s)$, $\|X(s)\|_{\mathrm{HS}}^2=2\|s\|^2$.
Thus $R_j=X(s_j)$ proves the first inequality in
\eqref{q444:twocross:residual-lower-bound}.

The exact expansion
\begin{align*}
 &(u+2v_0+t/24+2\omega^2)(v_0+5t/12)\\
 &=5t\omega^2/6
   +v_0(u+2v_0+t/24+2\omega^2)
   +(5t/12)(u+2v_0+t/24)
\end{align*}
has nonnegative remainder terms, the last strictly positive for $t>0$.
Together with $d_M\geq0$ and $e^2>0$ it gives
$\mathcal P>5e^2t/6$, proving the strict second inequality.
Finally $e^2=2+\operatorname{tr}M\geq2$,
$0\leq v_0\leq\operatorname{tr}M\leq L$, and $|m|\geq\delta$
give \eqref{q444:twocross:uniform-bound}. Since unitary conjugation
preserves the Hilbert--Schmidt norm, the same inequalities hold for
the complete spatial residuals in the original literal matter frame.
\end{proof}


\subsubsection{Scope, provenance, and exact replay}
\label{q444:twocross:scope-replay}

Theorem~\ref{q444:twocross:thm:exclusion} and Proposition~\ref{q444:twocross:prop:exceptional-inclusion}
together prove an empty admissible fibre at every point of the explicit
eight-dimensional two-cross linear subspace inside the full exceptional set.
The result does not prove global emptiness or produce a full solution.
It does not assert a result for arbitrary time-dependent gauge connections,
and it makes no quantum mass-gap assertion. The earlier cross--cross--block exclusion in
Theorem~\ref{q444:cx:thm:block} instead assumes $C,Z_1\in\mathfrak p$, $Z_2\in\mathfrak k$;
here only $Z_1,Z_2\in\mathfrak p$ were assumed and $C\in\mathfrak k$ was proved.

The retained literal system is
\eqref{q444:cx:eq:phi}--\eqref{q444:cx:eq:sp2}; its complete scalar
coordinates are given in \eqref{q444:complex-eight-coordinates} and
Proposition~\ref{q444:complex-full-scalar-equivalence}.
The ambient projection statement and determinant convention are
Theorem~\ref{q444:generic-projection-obstruction} and
\eqref{q444:generic-sylvester-map}. The complete source proof, exact
checker, and replay receipt for this additional exclusion are retained
under \path{supporting_materials/workbench/evidence/two_cross_exceptional_20260905/}.
All matrix identities needed for the theorem are derived above; no
external classification theorem, numerical root search, or unproved
literature identity is used.

The companion \path{check_two_cross_exceptional.py} checks 26 exact groups:
the original force and constants, both frame maps (the second over the
entire coordinate ring $\gamma\bar\gamma=1$), every scalar residual before
and after temporal classification, both inertia maps, the full bracket and
Gauss residual, both complete spatial residuals, the nonnegative determinant
identity, the full outer-product identity and the final positivity formula.
Every real component of $v,w$ is an independent symbol. It uses no floating
point acceptance criterion and writes no files. One serial run passed all
26 groups in $0.416$ seconds of measured checker time.
The determinant block-factor proof and all positivity/domain arguments
are analytic proofs in the preceding sections, not claims inferred from
sampled determinant values.

The additional denominator-free residual identity and the full spatial
error bound are proved in Proposition~\ref{q444:twocross:quantitative}.
The companion \path{check_residual_certificate.py} verifies its polynomial
identity and exact positive remainder. The Cauchy--Schwarz estimates and
their scalar--Gauss domain are proved in the text, not inferred from tests.
